\documentclass[11pt,a4paper]{article}
\usepackage[margin=2.5cm]{geometry}
\usepackage{amsmath,amssymb,amsthm}
\usepackage[T1]{fontenc}
\usepackage{lmodern}

\newtheorem{definition}{Definition}
\newtheorem{axiom}{Axiom}
\newtheorem{proposition}{Proposition}
\newtheorem{corollary}{Corollary}
\newtheorem{remark}{Remark}

\newcommand{\PhiPlus}{\Phi^{+}}
\newcommand{\card}[1]{\left|#1\right|}
\newcommand{\Real}{\mathbb{R}}

\title{A Mathematical Specification of the Macro-Philosophical Index $\Phi^{+}$}
\author{Bagirov Hasan}
\date{October 2026}

\begin{document}
\maketitle

\section{Introduction and scope}

This document gives a mathematical specification of the index $\PhiPlus$, the
Macro-Philosophical Index. The specification is purely theoretical. It fixes
scales, the aggregation formula, and a graph-theoretic construction for graded
logical coherence. It contains no empirical application, no calibration, and no
derivation of the numerical parameters: weights and the filter exponent are
declared parameters of a family of indices.

Let $\mathcal{T}$ be a class of formalized theories. For $T\in\mathcal{T}$ we
write $s_X(T)\in[0,1]$ for the score of $T$ on variable $X$, where
\[
X\in\{O,M,F,N,R\}.
\]
The value $0$ means complete absence of the corresponding property, and $1$
means maximal expression.

\section{Axioms}

\begin{axiom}[Continuous scale]
Every variable is measured on the continuous interval $[0,1]$:
\[
s_X(T)\in[0,1]\qquad\text{for all }X\in\{O,M,F,N,R\}.
\]
\end{axiom}

\begin{axiom}[No compensation]
\label{ax:no-comp}
For every $X\in\{O,M,F,N,R\}$ and every assignment of the remaining scores,
\[
s_X(T)=0\ \Longrightarrow\ \PhiPlus(T)=0.
\]
In particular, a zero on one variable cannot be offset by any values of the
other variables.
\end{axiom}

\begin{axiom}[Coherence filter]
\label{ax:filter}
Logical coherence $N$ enters multiplicatively with a positive exponent: there
exist $\lambda>0$ and a function $g\colon[0,1]^{4}\to[0,\infty)$ such that
\[
\PhiPlus(T)=g\bigl(s_O(T),s_M(T),s_F(T),s_R(T)\bigr)\,\bigl(s_N(T)\bigr)^{\lambda},
\]
where $g$ is nondecreasing in each argument and $g(1,1,1,1)=1$.
\end{axiom}

\begin{remark}
Axiom~\ref{ax:no-comp} is the precise content of the informal requirement that
the index not be a linear weighted sum. A linear sum violates it: with
$\sum_i \tilde w_i=1$ and $\tilde w_i>0$, one has
$\sum_i\tilde w_i s_i=1-\tilde w_k$ when $s_k=0$ and all other scores equal $1$.
\end{remark}

\section{Scales and measurement model}

\begin{definition}[Formalized theory]
A formalized theory $T$ consists of a finite set $V(T)$ of atomic theses and a
consequence relation $\vdash_T$.
\end{definition}

\begin{definition}[Measurement model]
For $X\in\{O,M,F,R\}$ let $C_X$ be a finite nonempty set of criteria and let
$\alpha_c>0$ $(c\in C_X)$ be criterion weights. A \emph{measurement model} for
$T$ is a family of degrees
\[
x_c^{X}(T)\in[0,1]\qquad(c\in C_X,\ X\in\{O,M,F,R\}),
\]
where $x_c^{X}(T)$ expresses how strongly $T$ satisfies criterion $c$ for
variable $X$. The specification fixes the sets $C_X$ and the weights
$\alpha_c$; the degrees $x_c^{X}(T)$ are supplied by the measurement model.
\end{definition}

\begin{definition}[Criterion-based variable score]
For $X\in\{O,M,F,R\}$,
\[
s_X(T)=\frac{\sum_{c\in C_X}\alpha_c\,x_c^{X}(T)}{\sum_{c\in C_X}\alpha_c}.
\]
Unless otherwise stated, $\alpha_c=1$ for all criteria, so that
\[
s_X(T)=\frac{1}{\card{C_X}}\sum_{c\in C_X}x_c^{X}(T).
\]
\end{definition}

\begin{definition}[Ontological Clarity $O$]
$C_O=\{O_1,O_2,O_3,O_4\}$, where
\begin{itemize}
\item $O_1$: $T$ explicitly fixes its existential commitments;
\item $O_2$: $T$ specifies the domain of quantification;
\item $O_3$: $T$ specifies identity or individuation conditions for primitive entities;
\item $O_4$: $T$ distinguishes existence claims from mere predication or modal status.
\end{itemize}
Then $s_O(T)=\frac14\sum_{j=1}^{4}x_{O_j}^{O}(T)$.
\end{definition}

\begin{definition}[Modal Differentiation $M$]
$C_M=\{M_1,M_2,M_3,M_4\}$, where
\begin{itemize}
\item $M_1$: $T$ distinguishes eternal from temporal entities or truths;
\item $M_2$: $T$ distinguishes immutable from mutable entities or truths;
\item $M_3$: $T$ distinguishes aseitic from dependent entities;
\item $M_4$: $T$ assigns modal statuses (necessity, contingency, possibility) to fundamental claims.
\end{itemize}
Then $s_M(T)=\frac14\sum_{j=1}^{4}x_{M_j}^{M}(T)$.
\end{definition}

\begin{definition}[Fundamentality $F$]
$C_F=\{F_1,F_2,F_3,F_4\}$, where
\begin{itemize}
\item $F_1$: $T$ identifies a fundamental level or fundamental entities;
\item $F_2$: $T$ specifies a grounding or dependence direction;
\item $F_3$: $T$ distinguishes non-fundamental from fundamental entities;
\item $F_4$: $T$ provides an explicit fundamentality predicate or ordering.
\end{itemize}
Then $s_F(T)=\frac14\sum_{j=1}^{4}x_{F_j}^{F}(T)$.
\end{definition}

\begin{definition}[Relational Completeness $R$]
$C_R=\{R_1,R_2,R_3,R_4\}$, where
\begin{itemize}
\item $R_1$: $T$ specifies primitive relations;
\item $R_2$: $T$ specifies domains and ranges of relations;
\item $R_3$: $T$ specifies composition or closure rules for relations;
\item $R_4$: $T$ states completeness of its relational inventory.
\end{itemize}
Then $s_R(T)=\frac14\sum_{j=1}^{4}x_{R_j}^{R}(T)$.
\end{definition}

The coherence score $s_N(T)$ is defined in Section~\ref{sec:coherence}.

\section{Aggregation}

\begin{definition}[Weight budget and exponents]
The nominal weights are
\[
w_O=0.30,\qquad w_M=0.20,\qquad w_F=0.20,\qquad w_R=0.15,\qquad w_N=0.15,
\]
with $\sum_{X}w_X=1$. Set
\[
W=w_O+w_M+w_F+w_R=1-w_N=0.85,
\qquad
\widetilde w_i=\frac{w_i}{W}\quad(i\in\{O,M,F,R\}),
\]
so that $\widetilde w_O+\widetilde w_M+\widetilde w_F+\widetilde w_R=1$. The
filter exponent $\lambda>0$ is an independent parameter; the default value is
$\lambda=0.5$.
\end{definition}

\begin{definition}[Macro-Philosophical Index $\PhiPlus$]
For $T\in\mathcal{T}$ and $\lambda>0$,
\[
\PhiPlus(T)
=
\Biggl(\prod_{i\in\{O,M,F,R\}} s_i(T)^{w_i}\Biggr)
\bigl(s_N(T)\bigr)^{\lambda}
=
\Biggl(\prod_{i\in\{O,M,F,R\}} s_i(T)^{\widetilde w_i}\Biggr)^{W}
\bigl(s_N(T)\bigr)^{\lambda},
\]
with the convention $0^{a}=0$ for $a>0$.
\end{definition}

\begin{remark}[Status of $w_N$]
\label{rem:wn}
The value of $\PhiPlus$ depends only on the parameters
$(w_O,w_M,w_F,w_R,\lambda)$. The nominal share $w_N$ enters no formula except
through the accounting identity $W=1-w_N$. The budget-consistent choice
$\lambda=w_N$ is available, but is not adopted here, so that filter intensity
can be varied independently of the budget.
\end{remark}

\begin{remark}[Not a weighted mean]
Since $W+\lambda=1.35\neq 1$ for the default parameters, $\PhiPlus$ is not a
weighted geometric mean of the five scores; its maximum $1$ is attained only at
the corner $(1,1,1,1,1)$, and it penalizes deficiency more strongly than a mean
with the same weights would.
\end{remark}

\begin{proposition}[Admissibility]
\label{prop:admissible}
The construction of $\PhiPlus$ satisfies Axioms~1--3.
\end{proposition}

\begin{proof}
Axiom~1 holds by definition of the scales and of $s_N\in[0,1]$
(Section~\ref{sec:coherence}). For Axiom~\ref{ax:no-comp}: if $s_X=0$ with
$X\in\{O,M,F,R\}$ then the corresponding factor $s_X^{w_X}$ vanishes, and since
$w_X>0$ the product vanishes; if $s_N=0$ then $s_N^{\lambda}=0$ because
$\lambda>0$. Axiom~\ref{ax:filter} holds with $g$ the weighted geometric mean of
$s_O,s_M,s_F,s_R$, which is nondecreasing and equals $1$ at $(1,1,1,1)$.
\end{proof}

\section{Coherence}
\label{sec:coherence}

\begin{definition}[Dependency graph]
For a formalized theory $T$, let $D(T)=(V,E)$ be the directed dependency graph
with vertex set $V=V(T)$ and edge set $E\subseteq V\times V$. An edge $(v,u)\in E$
means that thesis $v$ presupposes or depends on thesis $u$ for its justification.
\end{definition}

\begin{definition}[Feedback vertex set]
A set $H\subseteq V$ is a feedback vertex set of $D(T)$ if the induced subgraph
$D(T)[V\setminus H]$ is acyclic. Define
\[
h^{*}(T)=\min\bigl\{\card{H}: H\subseteq V,\ D(T)[V\setminus H]\text{ is acyclic}\bigr\},
\]
and let $H^{*}(T)$ be any minimizer. Equivalently, $H^{*}(T)$ is a minimum
hitting set of the family of directed cycles of $D(T)$.
\end{definition}

\begin{definition}[Graded coherence $s_N$]
Let $\chi(T)\in\{0,1\}$ be the consistency indicator,
\[
\chi(T)=
\begin{cases}
0, & \text{if } T\vdash_T\bot,\\
1, & \text{otherwise}.
\end{cases}
\]
Define
\[
s_N(T)=\chi(T)\left(1-\frac{h^{*}(T)}{\card{V(T)}}\right),
\]
with the convention $h^{*}(T)/\card{V(T)}:=0$ when $V(T)=\emptyset$.
\end{definition}

\begin{remark}[Two failure modes]
The score $s_N$ is a product of two components: the consistency indicator
$\chi$, which detects derivability of $\bot$ through $\vdash_T$, and the
\emph{acyclicity degree} $a(T)=1-h^{*}(T)/\card{V(T)}$, which penalizes circular
justification. Both must hold for $s_N$ to be positive. Consistency and
acyclicity are logically independent: a consistent theory may have circular
dependencies, and an inconsistent one need not.
\end{remark}

\begin{remark}[Size dependence]
\label{rem:size}
For fixed $h^{*}$, the acyclicity degree increases with $\card{V}$: a single
$3$-cycle in a graph with $100$ vertices gives $a=0.97$, whereas a loop on a
one-vertex graph gives $a=0$. The measure is therefore sensitive to theory size,
not only to the share of cyclic structure. This is accepted in the present
version; calibration will decide whether a size-corrected variant is preferable.
\end{remark}

\begin{definition}[Exact algorithm]
The minimum feedback vertex set is computed by the integer linear program
\[
\begin{aligned}
\min\quad & \sum_{v\in V} y_v\\
\text{subject to}\quad & \sum_{v\in C} y_v\ge 1 && \text{for every directed cycle } C \text{ of } D(T),\\
& y_v\in\{0,1\} && \text{for every } v\in V.
\end{aligned}
\]
If $y^{*}$ is optimal, then $H^{*}(T)=\{v\in V: y_v^{*}=1\}$ and
$h^{*}(T)=\sum_v y_v^{*}$.
\end{definition}

\begin{remark}[Complexity and implementation]
Feedback vertex set is NP-hard in general, so exact computation is exponential
in the worst case. In practice the cycle family can be generated lazily: solve
the relaxation on a current subset of cycles, find a violated cycle in a
strongly connected component of the resulting subgraph, add it as a cutting
plane, and iterate until optimality. The construction is exact; only its cost
depends on $\card{V}$.
\end{remark}

\section{Properties}

\begin{proposition}[Zero set]
\label{prop:zero}
$\PhiPlus(T)=0$ if and only if $s_N(T)=0$ or $s_i(T)=0$ for some
$i\in\{O,M,F,R\}$.
\end{proposition}

\begin{proof}
All factors are nonnegative and all exponents are positive. A product of
nonnegative reals with positive exponents vanishes exactly when one of its
factors vanishes; conversely, if all five scores are positive then all factors
are positive.
\end{proof}

\begin{proposition}[Monotonicity, bounds, continuity]
\label{prop:mono}
$\PhiPlus$ is continuous on $[0,1]^{5}$, takes values in $[0,1]$, equals $1$
only at $(1,1,1,1,1)$, and is nondecreasing in each argument, strictly so
whenever $\PhiPlus>0$.
\end{proposition}

\begin{proof}
Continuity: $x\mapsto x^{a}$ is continuous on $[0,1]$ for $a>0$ with $0^{a}=0$,
and a finite product of continuous functions is continuous. Bounds: each factor
lies in $[0,1]$, hence so does the product; the product equals $1$ only if every
factor equals $1$. Monotonicity: for $i\in\{O,M,F,R\}$ and $s_i>0$,
\[
\frac{\partial\PhiPlus}{\partial s_i}=\frac{w_i}{s_i}\,\PhiPlus\ge 0,
\]
with equality only when $\PhiPlus=0$; for $N$,
$\partial\PhiPlus/\partial s_N=(\lambda/s_N)\PhiPlus\ge 0$.
\end{proof}

\begin{proposition}[Anti-compensation]
\label{prop:anti}
Let $P(T)=\prod_{i\in\{O,M,F,R\}}s_i(T)^{w_i}$ with $w_i>0$. If $s_k(T)\to 0$
for some $k$, then $P(T)\to 0$ regardless of the other factors. Hence no finite
increase in the other variables can compensate a sufficiently low value of
$s_k$.
\end{proposition}

\begin{proof}
For $k\in\{O,M,F,R\}$,
\[
P(T)=s_k(T)^{w_k}\prod_{i\neq k}s_i(T)^{w_i}\le s_k(T)^{w_k},
\]
since $s_i(T)\le 1$ and $w_i>0$. If $s_k(T)\to 0$ then $s_k(T)^{w_k}\to 0$,
hence $P(T)\to 0$. The same bound holds for the full index, because
$s_N^{\lambda}\le 1$ when $s_N\le 1$.
\end{proof}

\begin{corollary}[Threshold dominance]
For every $\varepsilon>0$ and every $k\in\{O,M,F,R\}$,
\[
s_k(T)<\varepsilon^{1/w_k}\ \Longrightarrow\ \PhiPlus(T)<\varepsilon,
\]
independently of the values of the other four scores.
\end{corollary}

\begin{proof}
Immediate from $\PhiPlus(T)\le s_k(T)^{w_k}$, established in the proof of
Proposition~\ref{prop:anti}.
\end{proof}

\section{Status of operationalization}
\label{sec:status}

The specification operates on four distinct levels, which should not be
conflated.

\begin{enumerate}
\item \textbf{Formal structure} (complete). Carriers, scales, the aggregation
operation, the zero set, monotonicity, and the coherence algorithm are fully
determined: given a vector of degrees, every quantity is computable.
\item \textbf{Indicator specification} (complete at the level of content). The
sets $C_O,C_M,C_F,C_R$ enumerate what is to be assessed. The criteria are
ostensive: they name properties, not decision procedures.
\item \textbf{Measurement protocol} (deliberately absent). Coding rules, anchors
for the values $0$ and $1$, aggregation of multiple coders, and reliability
requirements are not part of this document. The degrees $x_c^{X}(T)$ are
primitives of the measurement model.
\item \textbf{Calibration} (deliberately absent). The weights
$(w_O,w_M,w_F,w_R)$, the filter exponent $\lambda$, and the choice of
size correction in $s_N$ are free parameters of a family of indices. No
empirical derivation is claimed.
\end{enumerate}

Consequently the document specifies a \emph{family} of indices, parametrized by
a measurement model and by $(\tilde w,\lambda)$. This is intentional: the
theoretical version fixes the architecture and leaves measurement and
calibration to a later stage.

\section{Conclusion}

We have specified the index
\[
\PhiPlus(T)
=
\Biggl(\prod_{i\in\{O,M,F,R\}} s_i(T)^{w_i}\Biggr)
\bigl(s_N(T)\bigr)^{\lambda},
\qquad
\lambda=0.5,
\]
with core weights
\[
w_O=0.30,\quad w_M=0.20,\quad w_F=0.20,\quad w_R=0.15,
\]
nominal share $w_N=0.15$ (recorded for accounting only, see
Remark~\ref{rem:wn}), and $W=0.85$.

The variables $O,M,F,R$ are graded on $[0,1]$ by criterion-based scores. The
coherence variable $N$ is graded as the product of a consistency indicator and
an acyclicity degree derived from the minimum feedback vertex set of the
dependency graph; logical contradiction forces $s_N=0$. The weighted geometric
core satisfies the no-compensation axiom in the precise form of
Axiom~\ref{ax:no-comp}, and the multiplicative filter with $\lambda>0$ makes
$\PhiPlus$ vanish on every theory with $s_N=0$. The zero set, monotonicity,
bounds, continuity, and the anti-compensation bound are established in
Section~6. The construction is purely theoretical and contains no empirical
calibration.

\end{document}